\documentclass[12pt]{article}
\usepackage{amsmath,amssymb,amsfonts,amsthm}
\usepackage[margin=1in]{geometry}

\begin{document}

\begin{center}
{\Large\bfseries Chapter 9, Problem 24}\\[6pt]
Byron \& Fuller, \textit{Mathematics of Classical and Quantum Physics}
\end{center}

\bigskip

\textbf{Problem.} It can be shown that $f_0^2(q_1) = \frac{1}{q_1}e^{i\delta_0}\sin\delta_0$, where $\delta_0$ is the $s$-wave phase shift. Using Fig.~9.4, obtain a graph of $\delta_0(q_1)$ versus $q_1$. Are the results from the two different curves in Fig.~9.4 consistent? To what extent is $\delta_0$ determined uniquely?

\bigskip
\textbf{Solution.}

\subsection*{Relation between $f_0$ and $\delta_0$}

The $s$-wave scattering amplitude is parameterized by the phase shift $\delta_0$ as:
\[
f_0(q_1) = \frac{1}{q_1}\,e^{i\delta_0}\sin\delta_0 = \frac{1}{q_1}\,\frac{e^{2i\delta_0} - 1}{2i}.
\]

This gives:
\[
\mathrm{Re}\,f_0 = \frac{\sin\delta_0\cos\delta_0}{q_1} = \frac{\sin 2\delta_0}{2q_1}, \qquad \mathrm{Im}\,f_0 = \frac{\sin^2\delta_0}{q_1}.
\]

From these:
\[
|f_0|^2 = \frac{\sin^2\delta_0}{q_1^2}, \qquad \mathrm{Im}\,f_0 = q_1|f_0|^2 \quad (\text{optical theorem}).
\]

\subsection*{Extracting $\delta_0$ from Fig.~9.4}

From the real and imaginary parts of $F_0(p_1)$ plotted in Fig.~9.4 (for $mV_0/p^2\hbar^2 = 1.4$):

\textbf{Method 1 (from $\mathrm{Im}\,f_0$):}
\[
\delta_0 = \arcsin\!\left(\sqrt{q_1\,\mathrm{Im}\,f_0}\right).
\]

\textbf{Method 2 (from both parts):}
\[
\tan\delta_0 = \frac{\mathrm{Im}\,f_0}{\mathrm{Re}\,f_0} = \frac{\sin^2\delta_0}{\sin\delta_0\cos\delta_0} = \tan\delta_0.
\]
Actually, more usefully: $\delta_0 = \arctan\!\left(\frac{\mathrm{Im}\,f_0}{\mathrm{Re}\,f_0}\right)$, but this only gives $\delta_0$ modulo $\pi$.

\textbf{Method 3 (using both real and imaginary parts together):}
\[
\delta_0 = \frac{1}{2}\arctan\!\left(\frac{2\,\mathrm{Re}(q_1 f_0)}{1 - 2q_1\,\mathrm{Im}\,f_0}\right).
\]

From Fig.~9.4:
\begin{itemize}
\item At small $p_1$: $F_0$ is primarily real and negative (repulsive potential), so $\delta_0 < 0$ and $|\delta_0|$ is small.
\item As $p_1$ increases, $|\delta_0|$ grows; the imaginary part grows while the real part decreases in magnitude.
\item At large $p_1$: $\delta_0 \to 0$ (Born regime, weak scattering at high energy).
\end{itemize}

\subsection*{Consistency of the two curves}

Fig.~9.4 shows both the exact result and the first iterative (Born) approximation. For the exact result, $\delta_0$ extracted from Re and Im parts must be consistent (same $\delta_0$ from either). For the Born approximation, $\mathrm{Im}\,f_0^{(1)} = 0$ (since it's real), which would give $\delta_0 = 0$ or $\pi$---clearly inconsistent with unitarity. The exact curves, however, satisfy $\mathrm{Im}\,f_0 = q_1|f_0|^2$ and yield a consistent $\delta_0$.

\subsection*{Uniqueness of $\delta_0$}

The phase shift $\delta_0$ is determined only modulo $\pi$ from the scattering amplitude, since:
\[
e^{i\delta_0}\sin\delta_0 = e^{i(\delta_0 + n\pi)}\sin(\delta_0 + n\pi) \quad \text{for integer } n.
\]

Wait---actually $\sin(\delta_0 + \pi) = -\sin\delta_0$ and $e^{i(\delta_0+\pi)} = -e^{i\delta_0}$, so the product is unchanged. Therefore $\delta_0$ is determined modulo $\pi$.

By convention, $\delta_0$ is chosen to be a continuous function of $q_1$ with $\delta_0 \to 0$ as $q_1 \to \infty$ (the Levinson theorem may add $n\pi$ at $q_1 = 0$ depending on the number of bound states). For a repulsive Yukawa potential with no bound states, $\delta_0(0) = 0$ and $\delta_0$ remains in $(-\pi/2, 0)$ for all $q_1$.

\[
\boxed{\delta_0 \text{ is determined modulo } \pi; \text{ the physical choice is fixed by continuity and } \delta_0(\infty) = 0.}
\]

\hfill $\blacksquare$

\end{document}
